\documentclass[11pt,a4paper]{article}
\usepackage[utf8]{inputenc}
\usepackage[spanish]{babel}
\usepackage{amsmath,amsfonts,amssymb}
\usepackage{graphicx}
\usepackage{geometry}
\geometry{margin=2.5cm}

\title{\textbf{Guía 2: Identificación experimental mediante prueba escalón}\\Control Inteligente -- ING01343-ING278}
\author{Politécnico Colombiano Jaime Isaza Cadavid\\Profesor: Deimer Miranda Montoya, MSc.(c).}
\date{}

\begin{document}
\maketitle

\section{Introducción}
Identificación experimental en lazo abierto de la función de transferencia del motor DC aplicando escalones de PWM al 30\%, 45\% y 60\%.

\section{Modelos Candidatos}
\begin{itemize}
    \item \textbf{Primer Orden sin retardo (FOP):}
    \begin{equation}
    G_{FOP}(s) = \frac{K}{\tau s + 1}
    \end{equation}
    \item \textbf{Primer Orden con retardo aparente (FOPDT):}
    \begin{equation}
    G_{FOPDT}(s) = \frac{K e^{-\theta s}}{\tau s + 1}
    \end{equation}
\end{itemize}

\end{document}
